% SM Table: Redocking Validation
% Generated from V4/V5 validation content, adapted for V7
% Canonical V7 package refresh: 2026-08-12

\begin{table}[!htbp]
\centering
\caption{Redocking validation: reproduction of co-crystallized ligand poses. \textit{Redocking assesses whether AutoDock Vina can reproduce experimentally determined binding poses. RMSD < 2.0 \si{\angstrom} indicates successful pose reproduction. Methotrexate (MTX) failed due to grid center positioning near NADPH cofactor rather than the folate binding site, illustrating target-specific limitations.}}
\label{tab:sm_redocking_validation}
\small
\begin{tabular}{llccp{6cm}}
\toprule
\textbf{Target} & \textbf{PDB} & \textbf{Ligand} & \textbf{RMSD (\si{\angstrom})} & \textbf{Result} \\
\midrule
PfDHFR & 7F3Y & MTX (methotrexate) & $\sim$30 & Failed (grid centered on NADPH) \\
PfDHFR & 7F3Y & P218 (pyrimethamine) & 1.42 & Success \\
PfCRT & 6UKJ & Y01 (proxy ligand) & 1.78 & Success \\
PfClpP & 2F6I & Peptide fragment & 1.65 & Success \\
PfATP4 & 9N10 & ADP (cofactor) & 1.23 & Success \\
\midrule
\multicolumn{5}{l}{\textit{Success rate: 4/5 alignable ligands (80\%)}} \\
\bottomrule
\multicolumn{5}{p{14cm}}{
\noindent\textbf{Notes:} Redocking performed using AutoDock Vina with the same grid parameters as virtual screening. RMSD calculated as heavy-atom deviation between top-scoring docked pose and co-crystallized ligand. MTX redocking failed because the grid was centered on the NADPH cofactor binding site ($\sim$12 \si{\angstrom} from folate pocket), not the MTX binding site; this illustrates target-specific validation requirements. PfCRT Y01 is a membrane-cavity proxy ligand (not a validated inhibitor). Success criterion: RMSD < 2.0 \si{\angstrom}.}
\end{tabular}
\end{table}
